\documentclass[12pt]{article}
\usepackage{amsmath, amssymb, geometry, graphicx}
\geometry{a4paper, margin=1in}

\title{\textbf{Methodology Deep-Dive: The Fortress Architecture v4.0}}
\author{Abhishek Raj}
\date{March 2026}

\begin{document}
\maketitle

\section{Introduction to Tiered Decoupling Method}
The core architectural thesis of this dissertation is the "Fortress Architecture," a Tiered Decoupling Method designed to isolate hardware drivers from high-level mission heuristics. This mono-repo, mono-logic approach is divided into three distinct modules: The Body, The Nerves, and The Brain.

\subsection{The Body: Physicality and Simulation}
The physical execution layer relies on the PX4 Autopilot firmware interfaced with a Gazebo physics simulation. This layer is strictly agnostic to agricultural logic. It processes standard MAVLink trajectories and outputs raw IMU and odometry data.

\subsection{The Nerves: Communication and QoS Durability}
The critical integration layer utilizes MicroXRCE-DDS to bridge PX4 with the ROS 2 Humble ecosystem. A fundamental methodological innovation was resolving the DDS Quality of Service (QoS) durability mismatch. By implementing a \texttt{TRANSIENT\_LOCAL} policy on the sensor topics, the system ensures that high-frequency IMU data streams are persistently available to late-joining nodes, entirely mitigating the "IMU Updates: 0" data starvation error.

\subsection{The Brain: V5 Master Mission Commander}
The highest tier is the V5 Master Mission Commander, which dictates the autonomous state machine (BOOT $\rightarrow$ PREFLIGHT $\rightarrow$ ARM $\rightarrow$ TAKEOFF $\rightarrow$ SURVEY). Under this decoupling, the Brain can trigger complex behaviors (e.g., sprayer actuation) without interfering with the low-level flight controller.

\section{Mathematical Innovation}

\subsection{Path Planning: Dynamic $A^*$ Optimization}
We implemented a dynamic $A^*$ search algorithm optimized for field traversal. The algorithm calculates the most efficient path avoiding dynamic obstacles (such as tractors) while adhering to maize rows. The cost function is defined as:
\begin{equation}
    f(n) = g(n) + h(n)
\end{equation}
where $g(n)$ is the exact cost from the start node to node $n$, and $h(n)$ is the heuristic estimated cost from node $n$ to the target.

\subsection{Robust State Estimation: ES-EKF on $SO(3)$}
To combat GPS-denied environments in the Munger riverine corridors, the Indra-Eye system deploys an Error-State Kalman Filter (ES-EKF). Standard EKFs utilizing Euler angles suffer from gimbal lock during aggressive wind-rejection maneuvers. To ensure mathematical robustness, our filter operates strictly on the $SO(3)$ Lie group manifold.

The error state vector is defined as:
\begin{equation}
    \delta \mathbf{x} = [\delta \mathbf{p}, \delta \mathbf{v}, \delta \boldsymbol{\theta}, \delta \mathbf{a}_b, \delta \boldsymbol{\omega}_b]^\top
\end{equation}
Here, $\delta \mathbf{p}$ is the position error, $\delta \mathbf{v}$ is the velocity error, $\delta \boldsymbol{\theta}$ is the orientation error vector (represented in the tangent space $\mathfrak{so}(3)$), and $\delta \mathbf{a}_b, \delta \boldsymbol{\omega}_b$ track IMU biases. This formulation guarantees singularity-free attitude integration, serving as a reliable "Inner Ear" for the UAV when satellite navigation is compromised.

\end{document}
